\section{Related Work}

%Recent developments in uncertainty quantification have made notable progress. One ubiquitous approach is Bayesian Neural Networks(BNN) with inherent posterior distribution over the model parameters to estimate the posterior predictive uncertainty. Although the exact posterior inference is intractable for over-parameterized models, several approximation approaches are proposed to alleviate this burden, such as Monte-Carlo Dropout~\citep{gal2016dropout}, Markov chain Monte Carlo(MCMC)~\citep{springenberg2016bayesian}, Variational Inference~\citep{graves2011practical}. Non-Bayesian approaches aim to avoid the problems of approximation and computational complexity in the Bayesian method. LULA~\citep{kristiadi2021learnable} train additional hidden units building on layers of MAP-trained model to improve its uncertainty calibration with Laplace approximation.

%Deep Ensemble~\citep{lakshminarayanan2017simple} quantifies the predictive uncertainty by training an ensemble of networks, while some other works~\citep{lee2017training, corbiere2019addressing} estimate the uncertainty by directly learning more calibrated prediction confidence. More recently, a class of approaches have been proposed to use a neural network to explicitly parameterize a Dirichlet prior distribution over the categorical label distributions~\citep{malinin2018predictive, sensoy2018evidential, malinin2019reverse, nandy2020towards, charpentier2020posterior, joo2020being}.  These works show advantages in multiple uncertainty quantification applications, such as out of distribution data(OOD) detection, due to its capability in measuring different kinds of uncertainties. 

%The post-hoc uncertainty learning problem is still under-explored, and few works focus on tackling this problem. Meta-model-based approaches train an auxiliary model to oversee the performance of the pretrained base model. White-box~\citep{chen2019confidence} trains multiple linear classifiers on intermediate layers of the pretrained model using additional validation data, then feed the outputs into another regression model to output a confidence score for predicting whether or not the pretrained model is correct on the validation samples. Similarly, DEUP~\citep{jain2021deup} trains an error predictor to estimate the epistemic uncertainty of the pretrained model in terms of the difference between generalization error and aleatoric uncertainty. However, many of these methods require additional data samples that are either validation or out of distribution dataset to train or tune the hyper-parameter, which is infeasible when these data are not available. Moreover, they are often not flexible enough to quantify and distinguish different types of uncertainties, such as epistemic uncertainty and aleatoric uncertainty, which are known to be crucial in various learning applications~\citep{hullermeier2021aleatoric}.

Uncertainty Quantification methods can be broadly classified as \textit{intrinsic} or \textit{extrinsic} depending on how the uncertainties are obtained from the machine learning models. Intrinsic methods encompass models that inherently provide an uncertainty estimate along with its predictions. Some intrinsic methods such as neural networks with homoscedastic/heteroscedastic noise models \citep{wakefield2013bayesian} and quantile regression \cite{koenker1978regression} can only capture \textit{Data} (aleatoric) uncertainty. Many applications including out-of-distribution detection, requires capturing both \textit{Data} (aleatoric) and \textit{Model} (epistemic) accurately. Bayesian methods such as Bayesian neural networks (BNNs) \cite{neal2012bayesian,blundell15,welling2011bayesian} and Gaussian processes \cite{gpbook} and ensemble methods~\citep{lakshminarayanan2017simple} are well known examples of intrinsic methods that can 
quantify both uncertainties. However, Bayesian methods and ensembles can be quite expensive and require several approximations to learn/optimize in practice~\citep{mackay1992practical,kristiadi2021learnable,welling2011bayesian}. Other approaches attempt to alleviate these issues by directly parameterizing a Dirichlet prior distribution over the categorical label proportions via a neural network~\citep{malinin2018predictive, sensoy2018evidential, malinin2019reverse, nandy2020towards, charpentier2020posterior, joo2020being}. 

Under model misspecification, Bayesian approaches are not well-calibrated and can produce severely miscalibrated uncertainties. In the particular case of BNNs, sparsity-promoting priors have been shown to produce better-calibrated uncertainties, especially in the small data regime \cite{ghosh2019model}, and somewhat alleviate the issue. Improved approximate inference methods and methods for prior elicitation in BNN models are active areas of research for forcing BNNs to produce better-calibrated uncertainties. Frequentist methods that approximate the jackknife have also been proposed to construct calibrated confidence intervals~\cite{alaa2019discriminative}.
%Open research questions along this direction include dealing with models with high dimensional parameters (modern neural networks), efficiently computing higher order Taylor series expansions likely necessary for UQ.


%Energy based methods have showed to more useful for OOD detection compared to the softmax scores \cite{liu2020energy}. The negative \textit{logsumexp} of the logits of deep neural networks can be defined as the energy function for samples $\mathbf{x}$ which fares much better in detecting OOD that the softmax scores. By using this energy and contrastive fine-tuning using auxiliary OOD dataset superior performance was obtained. Based on this definition of the energy function, maximising the value of softmax of a class is equivalent to maximizing the sum of the maximum logit value and the energy function. This hurts OOD detection as low energy values are preferred for in-distribution samples, which corresponds to high $p(x)$ for the in-distribution samples.

For models without an inherent notion of uncertainty, extrinsic methods are employed to extract uncertainties in a post-hoc manner. The post-hoc UQ problem is still under-explored, and few works focus on tackling this problem. %Meta-model-based approaches train an auxiliary model to oversee the performance of the pretrained base model. White-box~\citep{chen2019confidence} trains multiple linear classifiers on intermediate layers of the pretrained model using additional validation data, then feed the outputs into another regression model to output a confidence score for predicting whether or not the pretrained model is correct on the validation samples. 
One such approach is to build auxiliary or meta-models, which have been used successfully to generate reliable confidence measures (in classification) \citep{chen2019confidence}, prediction intervals (in regression), and to predict performance metrics such as accuracy on unseen and unlabeled data~\citep{elder2020learning}. Similarly, DEUP~\citep{jain2021deup} trains an error predictor to estimate the epistemic uncertainty of the pretrained model in terms of the difference between generalization error and aleatoric uncertainty. LULA~\citep{kristiadi2021learnable} trains additional hidden units building on layers of MAP-trained model to improve its uncertainty calibration with Laplace approximation. However, many of these methods require additional data samples that are either validation or out of distribution dataset to train or tune the hyper-parameter, which is infeasible when these data are not available. Moreover, they are often not flexible enough to distinguish epistemic uncertainty and aleatoric uncertainty, which are known to be crucial in various learning applications~\citep{hullermeier2021aleatoric}. In contrast, our proposed method does not require additional training data or modifying the training procedure of the base model.  


%One such approach is to build auxiliary or meta-models (MM) and have been used successfully to generate reliable confidence measures (in classification) \citep{chen2019confidence}, prediction intervals (in regression) and to predict performance metrics such as accuracy on unseen and unlabeled data~\citep{elder2020learning}. There exist approaches that modify the training of single model to extract or improve uncertainty quality. These method implicitly either learn a distance measure to distinguish between in-distribution and OOD~\cite{van2020simple}; or estimate the density of in-distribution data~\cite{liu2020energy}. Others use external datasets for calibration and to improve pre-trained models. The multi-class cross-entropy loss can be replaced with one-vs-rest loss based on the distance from class specific centroid \cite{van2020simple}. Here, the epistemic or model uncertainty corresponds to the distance from each center. The method can still be susceptible to low uncertainty far from decision boundary, to mitigate this a Jacobian regularization is added to promote sensitivity to non-discriminate features which is opposite of smoothness aimed at generalization. Energy based methods have showed to more useful for OOD detection compared to the softmax scores \cite{liu2020energy}. The negative \textit{logsumexp} of the logits of deep neural networks can be defined as the energy function for samples $\mathbf{x}$ which fares much better in detecting OOD that the softmax scores. By using this energy and contrastive fine-tuning using auxiliary OOD dataset superior performance was obtained.
%In contrast, we do not require modifying the training procedure of the base-model or updating it. We build on the success of meta-models and augment them to capture the data and model uncertainty more effectively.


%\noindent \textbf{UQ Applications/ Out-of-distribution detection:} \textcolor{red}{Write about OE, etc} Out-of-distribution samples have also been used to equip models with uncertainty. One recent approach \cite{chan2020unlabelled}, uses the OOD samples to learn a auxiliary network to predict if a sample is in-domain or OOD. The auxiliary network is constrained only to use the variance of the prediction distribution of the main network. This will lead to meaningful uncertainties as high predictive variance are mapped to OOD and low predictive distribution are mapped to in-domain samples. SDE-Net is similar in its use of OOD samples for training an auxiliary network, but use a Brownian motion based method to obtain predictive uncertainties to overcome uncertainty quality limitations of other methods.

